% Target ~950 words. Rubric: System Design + System Evaluation.
% Top band: multiple prototyping techniques, JUSTIFIED. A2 lost marks here for
% not explaining why each technique was chosen - do not repeat that.

\section{Methods}

\subsection{Ideation and design}

\FILL{Brainstorming and ideation process, ported from the A2 report and the design decision log.
Keep it short; the marks here are for justification, not narration.}

\subsection{Prototyping techniques, and why each was used}

The project used three prototyping techniques in sequence, each chosen for what it could reveal at
that stage.

\begin{description}
  \item[Paper sketches.] Used first because they are fast to produce and easy to discard, which
  suits a stage where the concept itself was still moving. \FILL{what the sketches revealed that
  justified moving on}
  \item[Storyboard.] Used to test the sequence of events across a whole session rather than
  individual screens, because the system's difficulty is coordination over time, not any single
  interface. \FILL{what it revealed}
  \item[Low- then high-fidelity Figma prototypes.] Low fidelity to settle layout and flow cheaply;
  high fidelity because the evaluation needed participants to react to something that behaves like
  a real application, including transitions and live state. \FILL{what changed between them}
\end{description}

\VERIFY{The A2 feedback said the appropriateness of the prototyping techniques was not adequately
explained. Each bullet above must say \emph{why that technique and not another}, not just what was
made.}

\subsection{Design guidance used}

\FILL{Which heuristics and guidelines informed the design - the community heuristics card deck
from Week 5, Nielsen's heuristics, and the platform design guidelines. Cite them properly; A2 was
marked down for references being too few or arbitrary.}

\subsection{Evaluation design}

The Assignment 3 brief requires at least one querying technique and one observation technique, with
a minimum of ten participants, five per technique. The study used \textbf{think-aloud} as the
observation technique and a \textbf{questionnaire} as the querying technique, with ten participants
in total and no overlap between the two groups.

\paragraph{Why think-aloud.} Observation was needed to find where people hesitate, misread the
interface or abandon a task, and think-aloud surfaces the reasoning behind those moments as they
happen rather than reconstructing it afterwards.

\paragraph{Why a questionnaire.} The questionnaire was chosen to measure what observation cannot:
whether a person would be willing to join a session with strangers, whether they would speak on a
voice channel, and whether they trust the privacy controls. These are dispositions rather than
behaviours, and a watched session cannot establish them. It also yields a System Usability Scale
score, giving a comparable number alongside the qualitative findings. A questionnaire can be
completed remotely, which allowed both halves of the study to run in parallel.

\paragraph{Instrument.} The questionnaire comprised the ten standard SUS items, six custom items
covering purpose comprehension, willingness to join strangers, comfort with voice, radar clarity,
trust in privacy controls and motivation, and two open questions.

\paragraph{Participants.} \FILL{How many were recruited, how they were found, relevant background
(whether they had used any 3D scanning application, whether they study the subject), and the age
spread. Do not name anyone; use the codes P1--P5 and Q1--Q5.}

\paragraph{Setting and procedure.} \FILL{Where and when sessions were run, on what device, who
moderated, and how long a session took. State that participants were briefed that the other
participants in the prototype are staged, since the system is synchronous and the prototype cannot
provide real co-present users.}

\paragraph{Tasks.} Both groups attempted the same five tasks, summarised in
Table~\ref{tab:tasks}.

\begin{table}[htbp]
\centering
\caption{Evaluation tasks, start screen and success criteria.}
\label{tab:tasks}
\small
\begin{tabularx}{\textwidth}{@{}lXl@{}}
\toprule
\textbf{\#} & \textbf{Task as read to the participant} & \textbf{Success criterion} \\
\midrule
T1 & Join the capture session happening now at the landmark. & Reaches the rally screen \\
T2 & Find which station you have been given, then tell the host you are ready without speaking. & Names station; uses a one-tap reply \\
T3 & Respond to the host's instruction to move, then say which side still needs photos. & Uses quick reply; reads the radar \\
T4 & Find the finished model and share it without revealing the exact location. & Reaches share with location off \\
T5 & A participant is photographing strangers' faces. Do something about it. & Opens participant menu and flags \\
\bottomrule
\end{tabularx}
\end{table}

\paragraph{Analysis.} Observations and open answers were consolidated into a single findings table,
with each problem rated for severity and assigned a decision to fix, defer or reject. SUS responses
were scored in the standard way.

\paragraph{Ethics.} \FILL{Consent, audio recording and deletion, anonymity. One short paragraph.}
